\documentclass[a4paper,12pt]{article}
\usepackage[T2A]{fontenc}
\usepackage[utf8]{inputenc}
\usepackage[russian]{babel}
\usepackage{amsmath,amssymb,mathtools}
\usepackage{geometry}
\usepackage{microtype}
\usepackage{tikz}
\usetikzlibrary{arrows.meta,positioning,shapes.geometric,calc}
\usepackage{tabularx,booktabs,longtable}
\usepackage{enumitem}
\usepackage{hyperref}
\usepackage{parskip}
\usepackage{fancyhdr}
\usepackage{setspace}
\usepackage{xcolor}
\usepackage{icomma}

\geometry{left=18mm,right=18mm,top=18mm,bottom=20mm}
\onehalfspacing
\renewcommand{\familydefault}{\sfdefault}
\setlength{\parindent}{0pt}
\setlength{\headheight}{15pt}
\setlist[itemize]{leftmargin=6mm,itemsep=2pt,topsep=3pt}
\setlist[enumerate]{leftmargin=8mm,itemsep=5pt,topsep=4pt}

\definecolor{accent}{HTML}{176B87}
\definecolor{darkaccent}{HTML}{0B4F63}
\definecolor{textgray}{HTML}{2E3438}
\definecolor{softgray}{HTML}{F2F5F6}
\definecolor{linegray}{HTML}{B8C4C9}

\color{textgray}
\hypersetup{colorlinks=true,linkcolor=darkaccent,urlcolor=darkaccent}
\pagestyle{fancy}
\fancyhf{}
\fancyhead[L]{\small НОД, алгоритм Евклида и НОК}
\fancyhead[R]{\small 29.09.26}
\fancyfoot[C]{\small\thepage}
\renewcommand{\headrulewidth}{0.3pt}
\renewcommand{\footrulewidth}{0pt}

\newcommand{\nodr}{\mathop{\text{НОД}}\nolimits}
\newcommand{\nokr}{\mathop{\text{НОК}}\nolimits}
\newcommand{\key}[1]{\textcolor{darkaccent}{\textbf{#1}}}
\newcommand{\sectionrule}{\vspace{-2mm}\par\noindent\textcolor{linegray}{\rule{\linewidth}{0.5pt}}\par\vspace{1mm}}

\begin{document}

% ---------- TITLE ----------
\thispagestyle{empty}
\begin{center}
\vspace*{28mm}
{\Large\textcolor{darkaccent}{Математика}}\\[9mm]
{\Huge\bfseries Чек-лист}\\[5mm]
{\LARGE\bfseries НОД, алгоритм Евклида и НОК}\\[15mm]
{\large Надя Клименко}\\[4mm]
{\normalsize 29.09.26}\\[22mm]
\begin{minipage}{0.78\textwidth}
\centering
\large
Главная идея занятия: научиться быстро находить НОД с помощью алгоритма Евклида и применять его в дробях, задачах на разбиение и при вычислении НОК.
\end{minipage}
\vfill
{\small Лёвин Артём Александрович эксклюзивно для Нади Клименко}
\end{center}

\begin{tikzpicture}[remember picture,overlay]
\draw[accent,line width=1.2pt]
  ($(current page.south west)+(18mm,17mm)$)
  .. controls ($(current page.south west)+(60mm,30mm)$)
  and ($(current page.south east)+(-65mm,7mm)$)
  .. ($(current page.south east)+(-18mm,19mm)$);
\draw[accent!45,line width=0.7pt]
  ($(current page.south west)+(18mm,12mm)$)
  .. controls ($(current page.south west)+(82mm,2mm)$)
  and ($(current page.south east)+(-80mm,26mm)$)
  .. ($(current page.south east)+(-18mm,13mm)$);
\end{tikzpicture}
\clearpage

% ---------- CORE ----------
\section*{1. Что нужно помнить про НОД}
\sectionrule
\key{Наибольший общий делитель} двух натуральных чисел --- самое большое натуральное число, на которое оба числа делятся без остатка.

Например,
\[
\nodr(32,12)=4,
\]
потому что 4 делит и 32, и 12, а большего общего делителя нет.

\textbf{Где НОД появляется в задачах:}
\begin{itemize}
  \item нужно разбить или разрезать объекты на \key{максимально крупные одинаковые части};
  \item нужно составить \key{максимальное число одинаковых наборов};
  \item нужно сократить дробь сразу на \key{наибольший общий множитель};
  \item нужно проверить, являются ли числа взаимно простыми.
\end{itemize}

\section*{2. Алгоритм Евклида}
\sectionrule
Вместо разложения больших чисел на простые множители можно последовательно делить с остатком.

Для $a>b>0$:
\begin{enumerate}
  \item разделите $a$ на $b$ с остатком: $a=bq+r$;
  \item замените пару $(a,b)$ на $(b,r)$;
  \item повторяйте, пока остаток не станет равен нулю;
  \item \key{последний ненулевой остаток} и есть $\nodr(a,b)$.
\end{enumerate}

\textbf{Пример.} Найдём $\nodr(52,20)$:
\[
\begin{aligned}
52&=20\cdot2+12,\\
20&=12\cdot1+8,\\
12&=8\cdot1+4,\\
8&=4\cdot2+0.
\end{aligned}
\]
Следовательно, $\nodr(52,20)=4$.

\textbf{Самопроверка:} найденный НОД обязан делить оба исходных числа без остатка.

\section*{3. Взаимно простые числа}
\sectionrule
Два натуральных числа называются \key{взаимно простыми}, если
\[
\nodr(a,b)=1.
\]

\textbf{Пример 1.} Для 221 и 65 алгоритм Евклида даёт НОД 13, поэтому эти числа не взаимно простые.

\textbf{Пример 2.} $1000=37\cdot27+1$. Как только в цепочке появился остаток 1, итоговый НОД будет равен 1. Значит, 1000 и 37 взаимно простые.

\section*{4. Сокращение дробей через НОД}
\sectionrule
Чтобы сократить дробь максимально за один шаг, удобно найти НОД числителя и знаменателя:
\[
\frac{a}{b}=\frac{a:\nodr(a,b)}{b:\nodr(a,b)}.
\]

\textbf{Пример.}
\[
\nodr(84,126)=42,
\qquad
\frac{84}{126}=\frac{84:42}{126:42}=\frac{2}{3}.
\]

Если после сокращения числитель и знаменатель взаимно просты, дробь несократима.

\clearpage
% ---------- FLOWCHART ----------
\thispagestyle{empty}
\begin{center}
{\Large\bfseries\textcolor{darkaccent}{Алгоритм Евклида}}\\[8mm]
\end{center}

\begin{center}
\begin{tikzpicture}[
  node distance=11mm,
  startstop/.style={ellipse,draw=accent,line width=0.9pt,minimum width=40mm,minimum height=11mm,align=center,fill=softgray},
  process/.style={rounded corners=3pt,draw=accent,line width=0.9pt,minimum width=92mm,minimum height=13mm,text width=82mm,align=center},
  decision/.style={diamond,aspect=2.2,draw=accent,line width=0.9pt,minimum width=50mm,minimum height=17mm,text width=40mm,align=center,inner sep=1pt},
  arrow/.style={-{Stealth[length=3mm]},line width=0.9pt,color=darkaccent}
]
\node (start) [startstop] {Даны натуральные числа $a>b>0$};
\node (divide) [process,below=of start] {Разделите $a$ на $b$ с остатком: $a=bq+r$};
\node (zero) [decision,below=14mm of divide] {$r=0$?};
\node (replace) [process,below=14mm of zero] {Замените пару $(a,b)$ на $(b,r)$};
\node (finish) [startstop,right=18mm of zero] {$\nodr$ равен текущему $b$};

\draw[arrow] (start) -- (divide);
\draw[arrow] (divide) -- (zero);
\draw[arrow] (zero) -- node[right,font=\small]{нет} (replace);
\draw[arrow] (replace.west) -- ++(-18mm,0) |- (divide.west);
\draw[arrow] (zero.east) -- node[above,font=\small]{да} (finish.west);
\end{tikzpicture}
\end{center}
\clearpage

% ---------- APPLICATIONS ----------
\section*{5. Прикладная задача на НОД: максимальный квадрат}
\sectionrule
Прямоугольник $96\times72$ см нужно покрыть одинаковыми максимально крупными квадратами без обрезков.

Сторона такого квадрата:
\[
\nodr(96,72)=24\text{ см}.
\]
Количество квадратов:
\[
\frac{96}{24}\cdot\frac{72}{24}=4\cdot3=12.
\]

\key{Сигнал НОД:} слова «максимально крупные одинаковые», «разделить без остатка», «наибольшее число одинаковых наборов».

\section*{6. Как связаны НОД и НОК}
\sectionrule
Для положительных натуральных чисел $a$ и $b$:
\[
\nodr(a,b)\cdot\nokr(a,b)=a\cdot b.
\]
Отсюда
\[
\boxed{\nokr(a,b)=\frac{a\cdot b}{\nodr(a,b)}}.
\]

\textbf{Пример.} Автобусы приходят каждые 24 и 36 минут. Найдём момент следующего совпадения:
\[
\nodr(24,36)=12,
\qquad
\nokr(24,36)=\frac{24\cdot36}{12}=72.
\]
Они снова придут одновременно через 72 минуты.

\section*{7. Как выбрать: НОД или НОК?}
\sectionrule
\begin{tabularx}{\linewidth}{@{}>{\bfseries}p{0.21\linewidth}X X@{}}
\toprule
 & НОД & НОК \\
\midrule
Что ищем & максимально возможный общий размер или максимальное число одинаковых групп & минимальный общий момент, длину периода или первое совпадение \\
Типичные слова & разрезать, разложить поровну, одинаковые наборы, самый большой размер & снова одновременно, через сколько впервые, повторяется, общий цикл \\
Главный вопрос & «Что должно делить оба числа?» & «Какое наименьшее число кратно обоим?» \\
\bottomrule
\end{tabularx}

\vspace{4mm}
\textbf{Полезная стратегия:} если нужен НОК больших чисел, сначала найдите НОД алгоритмом Евклида, затем используйте формулу связи НОД и НОК.

\section*{8. Частые ошибки}
\sectionrule
\begin{itemize}
  \item В алгоритме Евклида на следующем шаге берут \key{остаток}, а не частное.
  \item НОД --- это \key{последний ненулевой остаток}, а не ноль.
  \item Если НОД равен 1, числа взаимно простые; это не означает, что сами числа обязаны быть простыми.
  \item В задаче с квадратами НОД даёт \key{сторону} квадрата. Для количества квадратов нужно сравнить площади или посчитать число квадратов по каждой стороне.
  \item В задаче на повторяющиеся события обычно нужен НОК, потому что ищется \key{первое общее кратное времени}.
  \item В формуле $\nokr(a,b)=ab/\nodr(a,b)$ сначала удобно сократить, а затем умножать: так вычисления проще.
\end{itemize}

\clearpage
% ---------- TRAINING ----------
\section*{Тренировочный блок --- Путь А}
\sectionrule
\textbf{10 заданий.} Решайте по нарастающей сложности. В заданиях на НОД используйте алгоритм Евклида.

\begin{enumerate}
  \item Найдите $\nodr(84,30)$.
  \item Найдите $\nodr(252,198)$.
  \item Найдите $\nodr(2026,748)$.
  \item Являются ли числа 1000 и 37 взаимно простыми? Обоснуйте ответ через НОД.
  \item Сократите дробь $\dfrac{391}{299}$ до несократимого вида.
  \item Сократите дробь $\dfrac{132}{198}$ до несократимого вида.
  \item Прямоугольную площадку размером $168\times126$ см нужно покрыть одинаковыми максимально крупными квадратными плитками без подрезки. Найдите сторону плитки и количество плиток.
  \item Из 84 красных и 126 белых тюльпанов составляют максимально возможное число одинаковых букетов, используя все цветы. Сколько получится букетов? Сколько красных и белых тюльпанов будет в каждом букете?
  \item Два световых сигнала повторяются через каждые 90 с и 126 с. Сейчас они сработали одновременно. Через сколько секунд они впервые снова сработают одновременно? Ответ также запишите в минутах и секундах.
  \item Два рулона ленты длиной 840 см и 660 см нужно разрезать без остатка на одинаковые куски максимально возможной длины. Найдите длину одного куска и общее количество получившихся кусков.
\end{enumerate}

\clearpage
% ---------- ANSWERS ----------
\section*{Ответы}
\sectionrule
\begin{longtable}{@{}>{\bfseries}p{14mm}p{0.78\linewidth}@{}}
\toprule
№ & Ответ \\
\midrule
\endfirsthead
\toprule
№ & Ответ \\
\midrule
\endhead
1 & $6$ \\
2 & $18$ \\
3 & $2$ \\
4 & Да, взаимно простые: $\nodr(1000,37)=1$. \\
5 & $\dfrac{17}{13}$ \\
6 & $\dfrac{2}{3}$ \\
7 & Сторона плитки $42$ см; всего $12$ плиток. \\
8 & $42$ букета; в каждом по 2 красных и 3 белых тюльпана. \\
9 & $630$ с $=10$ мин $30$ с. \\
10 & Длина куска $60$ см; всего $25$ кусков. \\
\bottomrule
\end{longtable}

\vfill
\begin{center}
\small\textcolor{darkaccent}{Лёвин Артём Александрович эксклюзивно для Нади Клименко}
\end{center}

\end{document}
